\section{Limitations and Future Work}
%This paper has two limitations. Firstly, FunQA's input is limited to raw videos, which restricts the model's ability to process and analyze the content effectively. Secondly, our experimental results reveal a relatively low overall score. Therefore, future work could involve performing additional annotation of the video, such as panoramic segmentation of images, to create a more extensive and comprehensive VideoQA dataset. Additionally, exploring the development of novel evaluation metrics is a promising research direction in the VideoQA field. Secondly,  exploring the development of novel evaluation metrics holds promise as a research avenue within the VideoQA field.\\
%In conclusion, as a new benchmark, FunQA paves the way for exploring interesting videos and delving into deeper human understanding. By tackling the aforementioned limitations, future research can further enhance the performance of models in video comprehension tasks.
This paper has two limitations. \textbf{1)} The current FunQA dataset primarily contains video-level data and annotations. There is potential for enhanced video reasoning through denser annotations, akin to PVSG~\cite{yang2023panoptic}, which might include detailed spatial, temporal, and object-level annotations. \textbf{2)} The initial annotations were made in Chinese and later translated into English. While GPT was used to refine and complete the Chinese text, ensuring comprehensiveness and correctness, differences due to cultural differences between the languages might still persist.\\
Looking ahead, we plan to enrich FunQA with more detailed and varied annotations. We aim to develop new metrics for a more accurate evaluation of models, particularly for open-ended questions. Our goal is to steer models towards deeper video reasoning. However, to ensure fair comparisons and prevent data leakage, it is advised that future research does not utilize the FunQA testing set.